\documentclass[12pt,a4paper]{report}

\usepackage[utf8]{inputenc}
\usepackage[T1]{fontenc}
\usepackage[french]{babel}
\usepackage{geometry}
\geometry{margin=2.5cm}
\usepackage{hyperref}
\usepackage{setspace}
\usepackage{enumitem}
\usepackage{array}
\usepackage{longtable}
\usepackage{booktabs}
\usepackage{listings}
\usepackage{xcolor}
\usepackage{float}
\usepackage{chngcntr}

\renewcommand{\tablename}{Tableau}
\renewcommand{\figurename}{Figure}

\onehalfspacing

\hypersetup{
  colorlinks=true,
  linkcolor=black,
  urlcolor=blue,
  pdftitle={API REST sur TTGO T-Display},
  pdfauthor={Léo Ponchon, Yasser Razzaki, Léo Faucon}
}

\lstdefinestyle{code}{
  language=C++,
  basicstyle=\ttfamily\small,
  keywordstyle=\color{blue!70!black},
  commentstyle=\color{gray!70!black},
  stringstyle=\color{green!40!black},
  showstringspaces=false,
  frame=tb,
  framerule=0.3pt,
  columns=fullflexible,
  tabsize=2,
  breaklines=true
}

\counterwithin{figure}{section}
\counterwithin{table}{section}

\renewcommand{\thesection}{\arabic{section}}
\renewcommand{\thesubsection}{\thesection.\arabic{subsection}}
\renewcommand{\thesubsubsection}{\thesubsection.\arabic{subsubsection}}

\begin{document}

\begin{titlepage}
  \centering
  \vspace*{2cm}

  {\Huge Projet IoT\par}
  \vspace{0.5cm}
  {\LARGE API REST sur un TTGO T-Display\par}
  \vspace{1.5cm}

  {\Large HAI912I -- Partie IoT\par}
  \vspace{0.5cm}
  {\large Faculté des Sciences\par}
  \vfill

  \begin{tabular}{l}
    Léo Ponchon -- 22009418 \\
    Yasser Razzaki \\
    Léo Faucon
  \end{tabular}

  \vspace{1cm}

  Enseignant : Alban Mancheron\\
  Date : \today
\end{titlepage}

\tableofcontents
\newpage

\section{Introduction}

Ce projet porte sur la réalisation d'une petite station de mesure basée sur un
TTGO T-Display, lui-même construit autour d'un ESP32. La carte mesure la
température et la luminosité, affiche ces informations sur son écran intégré et
commande une LED en fonction de la luminosité.

L'objectif principal est de rendre ces fonctions accessibles au moyen d'une
API REST. Un client connecté au même réseau peut donc consulter les mesures,
connaître l'état de la LED et modifier son fonctionnement. Les deux boutons de
la carte permettent également de parcourir les informations affichées sans
utiliser l'API.

Le système associe acquisition analogique, affichage local, commande lumineuse
et communication réseau dans une même réalisation embarquée. La séparation du
programme en modules permet de distinguer les capteurs, l'affichage, la LED et
le serveur HTTP.

\section{Objectifs et cahier des charges}

Le programme doit fournir les fonctions suivantes :

\begin{itemize}
  \item mesurer la température à l'aide d'une thermistance ;
  \item mesurer une luminosité relative avec une photo-résistance ;
  \item lister les deux capteurs et récupérer leurs valeurs ;
  \item allumer, éteindre ou inverser l'état d'une LED par l'API ;
  \item définir deux seuils de luminosité pour un fonctionnement automatique ;
  \item utiliser l'écran et les boutons de la TTGO comme interface locale.
\end{itemize}

La réalisation conserve une organisation lisible et reproductible grâce à la
description du câblage, des broches et des fonctions principales.

\section{Matériel et montage}

\subsection{Composants utilisés}

Le montage utilise une TTGO T-Display, une photo-résistance, une thermistance
de 10 k$\Omega$, des résistances fixes, une LED et une résistance de protection
de 560 $\Omega$.

La thermistance et la photo-résistance sont des composants résistifs. Elles ne
possèdent donc pas de sens de branchement. En revanche, la LED est polarisée :
son anode est reliée au GPIO par la résistance de 560 $\Omega$ et sa cathode
est reliée à la masse.

\subsection{Affectation des broches}

\begin{table}[H]
  \centering
  \caption{Broches utilisées par le programme}
  \label{tab:broches}
  \begin{tabular}{ll}
    \toprule
    Fonction & GPIO \\
    \midrule
    Photo-résistance & 32 \\
    Thermistance & 33 \\
    LED externe & 25 \\
    Bouton suivant & 35 \\
    Bouton précédent & 0 \\
    Rétroéclairage & 4 \\
    \bottomrule
  \end{tabular}
\end{table}

Les GPIO32 et GPIO33 ont été choisis pour les mesures analogiques. Ils
appartiennent à l'ADC1 de l'ESP32, ce qui évite le conflit habituel entre
l'ADC2 et le Wi-Fi. L'alimentation des capteurs provient du 3,3 V de la carte.
Le 5 V n'est jamais envoyé sur une entrée GPIO.

Par convention, les fils rouges ou orange sont utilisés pour le 3,3 V, les
fils noirs ou bleus pour la masse et les fils jaunes ou blancs pour les
signaux. Cette convention facilite la vérification du montage avant la mise
sous tension.

\subsection{Ponts diviseurs}

La thermistance est câblée comme suit :

\begin{center}
\texttt{3V3 -- thermistance -- GPIO33 -- résistance 10 kOhm -- GND}
\end{center}

Le point relié à GPIO33 est une tension intermédiaire. Elle dépend de la
résistance de la thermistance, qui varie avec la température. La
photo-résistance est câblée de la même manière, avec son point intermédiaire
relié à GPIO32. Dans ce cas, le programme utilise directement la valeur ADC
pour comparer des niveaux de luminosité ; il ne prétend pas calculer une
luminosité en lux.

\section{Acquisition et traitement des mesures}

\subsection{Conversion analogique-numérique}

L'ESP32 est configuré avec une résolution de 12 bits. Les valeurs renvoyées
par l'ADC sont donc comprises entre 0 et 4095. Pour limiter les petites
variations, chaque mesure est la moyenne de 16 lectures espacées de quelques
millisecondes.

\begin{lstlisting}[style=code,caption={Moyenne de plusieurs lectures ADC}]
int lireMoyenne(uint8_t broche) {
  long somme = 0;

  for (int i = 0; i < 16; i++) {
    somme += analogRead(broche);
    delay(2);
  }

  return somme / 16;
}
\end{lstlisting}

\subsection{Température}

La thermistance utilisée est une NTC de 10 k$\Omega$ caractérisée par un
coefficient $B$ égal à 3950. La résistance est d'abord déduite de la valeur
ADC :

\[
R_{NTC}=R_{fixe}\left(\frac{4095}{ADC}-1\right)
\]

La température est ensuite calculée avec l'équation Beta :

\[
\frac{1}{T}=\frac{1}{T_0}+\frac{1}{B}\ln\left(\frac{R_{NTC}}{R_0}\right)
\]

Les températures sont exprimées en degrés Celsius après conversion depuis les
Kelvins. Si la valeur ADC est 0 ou 4095, la lecture est considérée comme
invalide et l'API renvoie \texttt{null} pour la température.

\subsection{Commande automatique de la LED}

La LED est allumée lorsque la luminosité passe sous le seuil d'allumage et
éteinte lorsqu'elle dépasse le seuil d'extinction. Les valeurs par défaut sont
1500 et 1900. L'écart entre les deux seuils forme une hystérésis et évite que
la LED ne change d'état en permanence lorsque la mesure oscille autour d'une
seule valeur.

Le mode automatique peut être désactivé par une commande manuelle
\texttt{on}, \texttt{off} ou \texttt{toggle}. La commande \texttt{auto}
réactive le comportement automatique.

\section{API REST}

Le serveur HTTP écoute sur le port 80 de l'ESP32. Son adresse IP est affichée
dans le moniteur série après la connexion au Wi-Fi. Les routes de lecture
utilisent la méthode GET et les modifications utilisent la méthode PUT.

La route \path{/} sert une interface web directement depuis l'ESP32. Cette
page affiche la température, la luminosité, l'état de la LED, son mode et les
seuils automatiques. Elle utilise uniquement les routes REST du projet pour
actualiser les données et envoyer les commandes. La page est stockée en mémoire
flash avec \texttt{PROGMEM}, afin de ne pas occuper inutilement la mémoire vive.
Un rappel des routes REST est également visible en bas de la page d'accueil.

\begin{longtable}{>{\raggedright\arraybackslash}p{5.2cm}>{\centering\arraybackslash}p{1.8cm}>{\raggedright\arraybackslash}p{7cm}}
  \caption{Routes de l'API REST}\\
  \toprule
  Route & Méthode & Fonction \\
  \midrule
  \endfirsthead
  \toprule
  Route & Méthode & Fonction \\
  \midrule
  \endhead
  \path{/} & GET & Interface web de supervision et de commande. \\
  \path{/api/v1/info} & GET & Adresse IP, puissance du signal Wi-Fi et durée de fonctionnement. \\
  \path{/api/v1/sensors} & GET & Liste logique des capteurs et mesures courantes. \\
  \path{/api/v1/sensors/temperature} & GET & Température courante. \\
  \path{/api/v1/sensors/light} & GET & Valeur ADC et niveau de luminosité. \\
  \path{/api/v1/led} & GET & État, mode et seuils de la LED. \\
  \path{/api/v1/led} & PUT & Allume la LED avec le paramètre \texttt{state=on}. \\
  \path{/api/v1/led} & PUT & Éteint la LED avec le paramètre \texttt{state=off}. \\
  \path{/api/v1/led} & PUT & Inverse l'état avec le paramètre \texttt{state=toggle}. \\
  \path{/api/v1/led} & PUT & Réactive le mode automatique avec \texttt{state=auto}. \\
  \path{/api/v1/threshold} & PUT & Modifie les seuils de luminosité avec les paramètres \texttt{on} et \texttt{off}. \\
  \bottomrule
\end{longtable}

La route \texttt{/api/v1/sensors} renvoie par exemple :

\begin{lstlisting}[style=code,caption={Exemple de réponse de l'API des capteurs}]
{
  "sensors": ["temperature", "light"],
  "light_adc": 1840,
  "light_level": "moyen",
  "thermistor_adc": 2050,
  "temperature_c": 23.7
}
\end{lstlisting}

Les erreurs de requête sont renvoyées avec le code HTTP 400. Une route
inconnue renvoie le code 404. Cette distinction permet au client de ne pas
confondre une erreur de chemin avec une mauvaise valeur de paramètre.

\section{Écran et boutons}

L'écran est organisé en trois pages :

\begin{enumerate}
  \item la température ;
  \item la valeur brute de la photo-résistance et son niveau ;
  \item l'état de la LED et son mode de fonctionnement.
\end{enumerate}

Le bouton relié à GPIO35 passe à la page suivante et celui relié à GPIO0
revient à la page précédente. Un délai de 250 ms est utilisé pour éviter
qu'un appui ne soit compté plusieurs fois. Le GPIO0 étant aussi lié au mode de
démarrage de l'ESP32, son état au démarrage peut influencer le téléversement.

\section{Validation fonctionnelle}

La validation du prototype a été réalisée progressivement, en commençant par
l'écran puis en ajoutant les capteurs, la LED et la communication Wi-Fi. Cette
méthode a permis de vérifier chaque fonction avant son intégration dans le
programme complet.

\begin{table}[H]
  \centering
  \caption{Résultats de la validation fonctionnelle}
  \label{tab:validation}
  \begin{tabular}{>{\raggedright\arraybackslash}p{4cm}>{\raggedright\arraybackslash}p{10cm}}
    \toprule
    Fonction & Résultat observé \\
    \midrule
    Écran & Les trois pages affichent successivement la température, la luminosité et l'état de la LED. \\
    Boutons & Les boutons permettent de parcourir les pages dans les deux sens. \\
    Photo-résistance & La valeur ADC varie selon l'éclairage reçu par le capteur. \\
    Thermistance & La température affichée varie lorsque la thermistance est chauffée ou touchée. \\
    LED & La LED répond aux commandes manuelles et au mode automatique piloté par la luminosité. \\
    Wi-Fi & La carte se connecte au réseau local et affiche son adresse IP dans le moniteur série. \\
    API REST & Les routes de lecture des capteurs et de commande de la LED sont disponibles sur le serveur HTTP. \\
    Gestion des erreurs & Les paramètres invalides renvoient HTTP 400 et les routes inconnues renvoient HTTP 404. \\
    \bottomrule
  \end{tabular}
\end{table}

\section{Gestion du microcontrôleur}

\subsection{Erreurs et reconnexion}

La connexion Wi-Fi possède un délai maximal de 20 secondes au démarrage. Si
elle échoue, le programme continue de faire fonctionner l'écran et les
capteurs, puis tente une reconnexion toutes les dix secondes. Une température
invalide est représentée par une erreur sur l'écran et par la valeur JSON
\texttt{null}.

Les paramètres de seuil sont contrôlés avant modification. Ils doivent être
compris entre 0 et 4095 et le seuil d'extinction doit être strictement
supérieur au seuil d'allumage.

\subsection{Configuration}

Les paramètres du réseau Wi-Fi sont lus dans le fichier local
\texttt{secrets.h}. Le fichier \texttt{secrets.h.example} décrit la structure
de cette configuration sans contenir de données personnelles. La bibliothèque
TFT\_eSPI utilise le fichier \path{Setup25_TTGO_T_Display.h} pour adapter
le pilote à l'écran de la TTGO T-Display.

\subsection{Autonomie}

Le prototype ne met pas l'ESP32 en veille : le Wi-Fi, l'écran et le
rétroéclairage restent actifs. Cette solution est adaptée à une alimentation
par USB, mais elle n'est pas optimisée pour une batterie. Une version orientée
autonomie pourrait réduire le rétroéclairage, espacer les mesures et utiliser
les modes de veille de l'ESP32.

\subsection{Sécurité}

L'API n'utilise pas de clé d'accès et fonctionne en HTTP. Son utilisation est
donc limitée au réseau local du projet. Une mise en production demanderait un
service intermédiaire en HTTPS avec authentification, plutôt qu'une ouverture
directe du port 80 vers l'ESP32. Les identifiants Wi-Fi restent séparés du
rapport.

\subsection{Mise à jour}

La programmation actuelle s'effectue par USB avec Arduino IDE. Les capteurs
sont alimentés en 3,3 V, tension utilisée par le montage. La mise à jour OTA
ne fait pas partie de cette version du prototype.

\section{Limites et améliorations possibles}

La luminosité est donnée en valeur ADC et non en lux. Une conversion physique
nécessiterait la caractérisation de la photo-résistance et de son pont
diviseur. La température dépend également de la valeur réelle de la
thermistance, de sa constante B et de la précision de la résistance fixe.

Les seuils et le mode de la LED sont conservés uniquement en mémoire. Ils
reviennent donc à leurs valeurs par défaut après un redémarrage. Une
amélioration consisterait à les enregistrer dans la mémoire non volatile de
l'ESP32, après avoir vérifié le nombre de cycles d'écriture.

Enfin, le programme a été séparé en plusieurs fichiers pour rendre son
organisation plus claire. Le fichier Arduino principal contient le démarrage,
la boucle principale et la connexion Wi-Fi. Les fichiers
\texttt{Capteurs.cpp/.h} regroupent les mesures analogiques et le calcul de la
température, tandis que \texttt{Led.cpp/.h} contient la commande de la LED et
la gestion des seuils. L'écran et les boutons sont regroupés dans
\texttt{Affichage.cpp/.h}, l'API REST dans \texttt{Api.cpp/.h} et
l'interface web dans \texttt{WebApp.cpp/.h}. Enfin,
\texttt{Config.h} rassemble les broches et les constantes du montage.

Cette séparation reste simple, mais elle facilite la lecture, les tests par
fonction et les modifications ultérieures. Elle évite aussi de mélanger dans
les mêmes fonctions la lecture des capteurs, l'affichage et les requêtes HTTP.

\section{Conclusion}

Le projet permet de réaliser une station de mesure simple autour d'un ESP32.
Les valeurs de la température et de la luminosité sont acquises sur des
entrées analogiques, affichées localement et rendues disponibles au moyen
d'une API REST. La LED peut être commandée manuellement ou automatiquement à
l'aide de seuils configurables.

Le montage a été construit progressivement afin de tester chaque composant
avant d'ajouter la fonction suivante. Cette méthode a permis de préserver le
matériel et de repérer une erreur de composant lors du test de la
thermistance. Les principales limites du prototype sont l'absence de
chiffrement, la luminosité exprimée en valeur ADC et la perte des seuils lors
d'un redémarrage. Elles sont prises en compte dans les perspectives
d'amélioration.

\section*{Références}
\addcontentsline{toc}{section}{Références}

\begin{itemize}
  \item Sujet du projet fourni dans le cadre de l'UE HAI912I.
  \item Tutoriel de prise en main du TTGO T-Display fourni par l'enseignant.
  \item Documentation du kit Adafruit Parts Pal :
  \url{https://www.adafruit.com/product/2975}.
  \item Documentation de la bibliothèque TFT\_eSPI et de la configuration
  TTGO T-Display.
\end{itemize}

\end{document}
